% 摘要与关键词
% 合并时在 main.tex 导言区 % 摘要与关键词
% 合并时在 main.tex 导言区 % 摘要与关键词
% 合并时在 main.tex 导言区 % 摘要与关键词
% 合并时在 main.tex 导言区 \input{chapters/abstract}，或把下面两行粘进导言区。
% 摘要 778 汉字（按汉字计），关键词五个。
% A-01--A-08、A-15、A-22--A-25、A-27、A-28、A-38；B-01 摘要数字来源。

\reportabstract{%
\linespread{1}\selectfont
ML-KEM-768（FIPS 203）的解封装在 Fujisaki-Okamoto（FO）变换下依次执行解密、重加密、比较与隐式拒绝，四步都接触私钥或解密消息。完整解封装的耗时在五十微秒量级，真正与消息位相关的差异只有数十纳秒；单次调用不可分辨，必须靠每类一万次、三次独立运行与采集前冻结的检验家族，才能把差异从噪声里拉出来。源码写成无分支，不等于编译产物无分支，判定必须绑定具体的二进制。检测工具又各有分工：dudect 与 RTLF 只回答两类耗时是否可区分，Memcheck 与 Microwalk 只回答哪里依赖秘密，Binsec 只回答给定边界内是否常数时间。单工具结论不能互相替代，未检出也常被读成安全；没有检测预算的阴性，读者不知道它排除了多大的效应。这一步的判定止于两类分布是否可区分，不进入密钥恢复。
本作品面向 ML-KEM-768 便携 C 实现，构建检测—定位—防护—复测平台。被测对象分三类：M1 与 M2 是真实历史案例，M3 与 M4 是人工校准夹具，M5 是阴性对照。平台在采集前冻结主类对、样本量、停止条件与 Holm 家族，用锁频控制器固定工作点并核对退出恢复；统计层用 Welch 检验、Holm 校正、逆方差固定效应合并与最小可检测效应量化每个阴性结论；定位层用六类工具交叉，并以三实现、五步骤、六类来源的 90 格排查逐格登记。定位、计时与防护复测绑定到同一份编译产物，产物摘要、编译条件与反汇编随每次判定登记，失败尝试与未运行项单独登记。\par
\par
M1 未修版本的两类密文均值差为 32.3 ns，约占 57 $\mu$s 的 0.06\%，三次运行方向一致；官方修复后区间含 0，解封装开销增加 1.51\%。M2 在三处秘密除法上定位阳性，完整解封装与除法热点微基准都未检出，检测预算为 25 至 67 ns；定位阳性不等于计时阳性，修复前后也都不能说明泄漏降低。M4 的提前返回臂可区分、常数时间臂不可区分，结论只覆盖 32 字节比较夹具，不外推到完整解封装。M5 在本机、本预算下未检出，运行过的六种工具误报为 0。KyberSlash 官方修复版本仍残留与编译器引入分支同源的消息编码分支；同一份 M1 源码在 aarch64 上被编成无分支形态，机制不复现；同机网络各档在本预算下未检出，不推断真实远程。
结论限定在单台 x86\_64 裸金属、clang 18.1.3、本地白盒与不做密钥恢复的范围内；时间结论不跨架构、不跨环境合并，每个未执行单列，不并进阴性。题面 12 条按 10 条满足、2 条部分登记，缺口与影响写在附录 A。%
}

\reportkeywords{ML-KEM；时间侧信道；常数时间实现；泄漏检测与定位；预注册统计}
，或把下面两行粘进导言区。
% 摘要 778 汉字（按汉字计），关键词五个。
% A-01--A-08、A-15、A-22--A-25、A-27、A-28、A-38；B-01 摘要数字来源。

\reportabstract{%
\linespread{1}\selectfont
ML-KEM-768（FIPS 203）的解封装在 Fujisaki-Okamoto（FO）变换下依次执行解密、重加密、比较与隐式拒绝，四步都接触私钥或解密消息。完整解封装的耗时在五十微秒量级，真正与消息位相关的差异只有数十纳秒；单次调用不可分辨，必须靠每类一万次、三次独立运行与采集前冻结的检验家族，才能把差异从噪声里拉出来。源码写成无分支，不等于编译产物无分支，判定必须绑定具体的二进制。检测工具又各有分工：dudect 与 RTLF 只回答两类耗时是否可区分，Memcheck 与 Microwalk 只回答哪里依赖秘密，Binsec 只回答给定边界内是否常数时间。单工具结论不能互相替代，未检出也常被读成安全；没有检测预算的阴性，读者不知道它排除了多大的效应。这一步的判定止于两类分布是否可区分，不进入密钥恢复。
本作品面向 ML-KEM-768 便携 C 实现，构建检测—定位—防护—复测平台。被测对象分三类：M1 与 M2 是真实历史案例，M3 与 M4 是人工校准夹具，M5 是阴性对照。平台在采集前冻结主类对、样本量、停止条件与 Holm 家族，用锁频控制器固定工作点并核对退出恢复；统计层用 Welch 检验、Holm 校正、逆方差固定效应合并与最小可检测效应量化每个阴性结论；定位层用六类工具交叉，并以三实现、五步骤、六类来源的 90 格排查逐格登记。定位、计时与防护复测绑定到同一份编译产物，产物摘要、编译条件与反汇编随每次判定登记，失败尝试与未运行项单独登记。\par
\par
M1 未修版本的两类密文均值差为 32.3 ns，约占 57 $\mu$s 的 0.06\%，三次运行方向一致；官方修复后区间含 0，解封装开销增加 1.51\%。M2 在三处秘密除法上定位阳性，完整解封装与除法热点微基准都未检出，检测预算为 25 至 67 ns；定位阳性不等于计时阳性，修复前后也都不能说明泄漏降低。M4 的提前返回臂可区分、常数时间臂不可区分，结论只覆盖 32 字节比较夹具，不外推到完整解封装。M5 在本机、本预算下未检出，运行过的六种工具误报为 0。KyberSlash 官方修复版本仍残留与编译器引入分支同源的消息编码分支；同一份 M1 源码在 aarch64 上被编成无分支形态，机制不复现；同机网络各档在本预算下未检出，不推断真实远程。
结论限定在单台 x86\_64 裸金属、clang 18.1.3、本地白盒与不做密钥恢复的范围内；时间结论不跨架构、不跨环境合并，每个未执行单列，不并进阴性。题面 12 条按 10 条满足、2 条部分登记，缺口与影响写在附录 A。%
}

\reportkeywords{ML-KEM；时间侧信道；常数时间实现；泄漏检测与定位；预注册统计}
，或把下面两行粘进导言区。
% 摘要 778 汉字（按汉字计），关键词五个。
% A-01--A-08、A-15、A-22--A-25、A-27、A-28、A-38；B-01 摘要数字来源。

\reportabstract{%
\linespread{1}\selectfont
ML-KEM-768（FIPS 203）的解封装在 Fujisaki-Okamoto（FO）变换下依次执行解密、重加密、比较与隐式拒绝，四步都接触私钥或解密消息。完整解封装的耗时在五十微秒量级，真正与消息位相关的差异只有数十纳秒；单次调用不可分辨，必须靠每类一万次、三次独立运行与采集前冻结的检验家族，才能把差异从噪声里拉出来。源码写成无分支，不等于编译产物无分支，判定必须绑定具体的二进制。检测工具又各有分工：dudect 与 RTLF 只回答两类耗时是否可区分，Memcheck 与 Microwalk 只回答哪里依赖秘密，Binsec 只回答给定边界内是否常数时间。单工具结论不能互相替代，未检出也常被读成安全；没有检测预算的阴性，读者不知道它排除了多大的效应。这一步的判定止于两类分布是否可区分，不进入密钥恢复。
本作品面向 ML-KEM-768 便携 C 实现，构建检测—定位—防护—复测平台。被测对象分三类：M1 与 M2 是真实历史案例，M3 与 M4 是人工校准夹具，M5 是阴性对照。平台在采集前冻结主类对、样本量、停止条件与 Holm 家族，用锁频控制器固定工作点并核对退出恢复；统计层用 Welch 检验、Holm 校正、逆方差固定效应合并与最小可检测效应量化每个阴性结论；定位层用六类工具交叉，并以三实现、五步骤、六类来源的 90 格排查逐格登记。定位、计时与防护复测绑定到同一份编译产物，产物摘要、编译条件与反汇编随每次判定登记，失败尝试与未运行项单独登记。\par
\par
M1 未修版本的两类密文均值差为 32.3 ns，约占 57 $\mu$s 的 0.06\%，三次运行方向一致；官方修复后区间含 0，解封装开销增加 1.51\%。M2 在三处秘密除法上定位阳性，完整解封装与除法热点微基准都未检出，检测预算为 25 至 67 ns；定位阳性不等于计时阳性，修复前后也都不能说明泄漏降低。M4 的提前返回臂可区分、常数时间臂不可区分，结论只覆盖 32 字节比较夹具，不外推到完整解封装。M5 在本机、本预算下未检出，运行过的六种工具误报为 0。KyberSlash 官方修复版本仍残留与编译器引入分支同源的消息编码分支；同一份 M1 源码在 aarch64 上被编成无分支形态，机制不复现；同机网络各档在本预算下未检出，不推断真实远程。
结论限定在单台 x86\_64 裸金属、clang 18.1.3、本地白盒与不做密钥恢复的范围内；时间结论不跨架构、不跨环境合并，每个未执行单列，不并进阴性。题面 12 条按 10 条满足、2 条部分登记，缺口与影响写在附录 A。%
}

\reportkeywords{ML-KEM；时间侧信道；常数时间实现；泄漏检测与定位；预注册统计}
，或把下面两行粘进导言区。
% 摘要 778 汉字（按汉字计），关键词五个。
% A-01--A-08、A-15、A-22--A-25、A-27、A-28、A-38；B-01 摘要数字来源。

\reportabstract{%
\linespread{1}\selectfont
ML-KEM-768（FIPS 203）的解封装在 Fujisaki-Okamoto（FO）变换下依次执行解密、重加密、比较与隐式拒绝，四步都接触私钥或解密消息。完整解封装的耗时在五十微秒量级，真正与消息位相关的差异只有数十纳秒；单次调用不可分辨，必须靠每类一万次、三次独立运行与采集前冻结的检验家族，才能把差异从噪声里拉出来。源码写成无分支，不等于编译产物无分支，判定必须绑定具体的二进制。检测工具又各有分工：dudect 与 RTLF 只回答两类耗时是否可区分，Memcheck 与 Microwalk 只回答哪里依赖秘密，Binsec 只回答给定边界内是否常数时间。单工具结论不能互相替代，未检出也常被读成安全；没有检测预算的阴性，读者不知道它排除了多大的效应。这一步的判定止于两类分布是否可区分，不进入密钥恢复。
本作品面向 ML-KEM-768 便携 C 实现，构建检测—定位—防护—复测平台。被测对象分三类：M1 与 M2 是真实历史案例，M3 与 M4 是人工校准夹具，M5 是阴性对照。平台在采集前冻结主类对、样本量、停止条件与 Holm 家族，用锁频控制器固定工作点并核对退出恢复；统计层用 Welch 检验、Holm 校正、逆方差固定效应合并与最小可检测效应量化每个阴性结论；定位层用六类工具交叉，并以三实现、五步骤、六类来源的 90 格排查逐格登记。定位、计时与防护复测绑定到同一份编译产物，产物摘要、编译条件与反汇编随每次判定登记，失败尝试与未运行项单独登记。\par
\par
M1 未修版本的两类密文均值差为 32.3 ns，约占 57 $\mu$s 的 0.06\%，三次运行方向一致；官方修复后区间含 0，解封装开销增加 1.51\%。M2 在三处秘密除法上定位阳性，完整解封装与除法热点微基准都未检出，检测预算为 25 至 67 ns；定位阳性不等于计时阳性，修复前后也都不能说明泄漏降低。M4 的提前返回臂可区分、常数时间臂不可区分，结论只覆盖 32 字节比较夹具，不外推到完整解封装。M5 在本机、本预算下未检出，运行过的六种工具误报为 0。KyberSlash 官方修复版本仍残留与编译器引入分支同源的消息编码分支；同一份 M1 源码在 aarch64 上被编成无分支形态，机制不复现；同机网络各档在本预算下未检出，不推断真实远程。
结论限定在单台 x86\_64 裸金属、clang 18.1.3、本地白盒与不做密钥恢复的范围内；时间结论不跨架构、不跨环境合并，每个未执行单列，不并进阴性。题面 12 条按 10 条满足、2 条部分登记，缺口与影响写在附录 A。%
}

\reportkeywords{ML-KEM；时间侧信道；常数时间实现；泄漏检测与定位；预注册统计}
